\begin{lemma}[Finite execution of the colour-round induction]
Fix $d,a,C\in[0,\infty]$, $\Delta,n\in\mathbb N$, and put $Q=\Delta^2+1$.
Fix exposure sequences $p,\rho$ and round sequences
$\varepsilon,\omega,\theta,\zeta,\widehat\zeta$, all mapping $\mathbb N$
to $[0,\infty]$, and a certificate $\mathcal K$ of
\noderef{InductionApplicationCertificate} for
$(\Delta,Q,d,\theta,p,\varepsilon,\rho,\zeta,\omega,\widehat\zeta)$.
Assume
\[
 p_0=d,\qquad p_{x+1}=\rho_x(d/2)\ (x<\Delta),\qquad
 0<a\leq1,\quad C/n\leq1,\quad \Delta(1-a)^n\leq\tfrac12.
\]
For every round $s<Q$ assume all of the following inequalities:
\[
 \varepsilon_s\leq\varepsilon_{s+1}\leq1,\qquad
 2\sqrt{\zeta_s}\leq\theta_s,\quad
 \zeta_s+\Delta\sqrt{\zeta_s}\leq\theta_s,\quad
 \zeta_s+2\Delta\zeta_s\leq\theta_s,
\]
\[
 0<\theta_s,\quad4\theta_s<\rho_x,\quad
 2\zeta_s<\rho_x,\quad2\zeta_s<1-\rho_x\qquad(x\leq\Delta),
\]
\[
 2\zeta_s<d/2,\quad2\zeta_s<1-d/2,\quad
 p_{x+1}+\varepsilon_s+1/n\leq p_x\ (x<\Delta),\quad
 \omega_{s+1}<\varepsilon_{s+1}.
\]
Writing the certificate fields as $N^{\rm wit}_{s,x},N^{\rm host}_s,
N^{\rm mat}_{s,x},N^{\rm exist}_{s,x},M_{s,x},
N^{\rm FG}_{s,x,y},a_{s,x,y}$, assume
\[
 N^{\rm wit}_{s,x},N^{\rm host}_s,N^{\rm mat}_{s,x},
 N^{\rm exist}_{s,x}\leq n,\quad M_{s,x}\leq C/2\qquad(x\leq\Delta),
\]
\[
 N^{\rm FG}_{s,x,y}\leq n,\quad a\leq a_{s,x,y}
 \qquad(x,y\leq\Delta).
\]
Finally assume, for every $s<Q$,
\[
 n^{-1}\leq\widehat\zeta_s<\zeta_s,\quad
 \widehat\zeta_s+2/n\leq\zeta_s,\quad
 \widehat\zeta_s+1/(\zeta_sn)\leq\zeta_s,
\]
\[
 \omega_{s+1}+2/n\leq\varepsilon_{s+1},\qquad
 \omega_{s+1}+1/(\varepsilon_{s+1}n)\leq\varepsilon_{s+1}.
\]
Then the following holds for every finite blow-up instance $B$ with
candidacy system $A$ and reduced graph on $[r]$. Suppose
$|X_j|=|V_j|\in\{n,n+1\}$ for every $j$, $\Delta(R)\leq\Delta$,
\noderef{InstanceSuperRegular} and \noderef{CandidacySuperRegular} hold
with $(\varepsilon_0,d)$, and \noderef{SaturatingSourcePair} holds on
every reduced edge. There are a colouring $c:[r]\to\{0,\ldots,Q-1\}$
and a specified law $\mu$ on compatible partial embeddings with domain
all of $V_H$ such that
\[
 \mathcal I(Q-1,B_{Q-1};V_H,\mu),\qquad
 B_{Q-1}=\{j:c(j)=Q-1\},
\]
in the sense of \noderef{ColourRoundInvariant}, with
$(\varepsilon,p,\Delta,C,n)$. In particular the full-domain restricted
maps have the $(C/n)$-vertex-spread bound. The final round is $Q-1$,
not $Q$.
\end{lemma}
\begin{proof}
The corrected paper processes each finite colour class sequentially and
then changes only the round bookkeeping. We give that finite execution,
leaving all local graph and probability estimates to the stated local
transition.

By \noderef{ReducedSquareColoring}, choose $q>0$, $q\leq Q$, and a
colouring $c_0:[r]\to\operatorname{Fin}(q)$ separating distinct adjacent
vertices and distinct vertices with a common neighbour. Define
$c(j)=\operatorname{Fin.castLE}(q\leq Q)(c_0(j))$.
This lift does not change the natural-number value of any colour.
It is injective on colour values: equality after lifting gives equality
of the original values and hence, by extensionality of finite indices,
equality in $\operatorname{Fin}(q)$. Thus $c$ satisfies precisely the
square-colouring separation used below. It takes values in
$\operatorname{Fin}(Q)$; unused colours are allowed. For each reduced
vertex $j$, the finset $\{k:R.Adj(j,k)\}$ has cardinality
$\deg_R(j)\leq\Delta(R)\leq\Delta$. This supplies the local per-vertex
degree hypothesis; it is valid also when there are no reduced vertices.

We next verify the one-cluster transition for any legal round $s<Q$.
Apply \noderef{VariableSizeInductionStep} with the substitution
\[
 (\theta,\mathrm{density},\epsilon,\rho,\eta,\omega,\xi)
       =(\theta,p,\varepsilon,\rho,\zeta,\omega,\widehat\zeta),
 \qquad D=\Delta,\quad q=Q.
\]
Use exactly the given certificate $\mathcal K$. Its type is the
\noderef{InductionApplicationCertificate} type required by this
substitution, with independent exposure and round indices.
The round bound is $s<Q$; the first displayed round inequalities give
$\varepsilon_s\leq\varepsilon_{s+1}$ and $\varepsilon_{s+1}\leq1$.
The derived degree bounds and the lifted square colouring give the next
two hypotheses. The assumed source/host size equalities give both
alternatives $n,n+1$ and the balance condition required there.
The given \noderef{InstanceSuperRegular} hypothesis has exactly error
$\varepsilon_0$ and density $d$, and the assumed
\noderef{SaturatingSourcePair} supplies its source-matching hypothesis.

The assumptions $0<a\leq1$, $C/n\leq1$, and
$\Delta(1-a)^n\leq1/2$ give the three constant/size bounds.
The three square-root and linear loss bounds give the simultaneous
good-edge inequalities. For every $x\leq\Delta$, the next display
gives $0<\theta_s$, $4\theta_s<\rho_x$, and both density and complementary
density bounds for $2\zeta_s$; the two corresponding host bounds
are also explicitly assumed. The recursion supplies
$p_{x+1}=\rho_x(d/2)$ for every $x<\Delta$, and the assumed density-gap
inequality supplies $p_{x+1}+\varepsilon_s+1/n\leq p_x$.
The output comparison is $\omega_{s+1}<\varepsilon_{s+1}$.
All four scalar thinning/matching thresholds and the comparison
$M_{s,x}\leq C/2$ are given for every exposure $x\leq\Delta$.
For every ordered pair $x,y\leq\Delta$ the Four Graphs threshold and
$a\leq a_{s,x,y}$ are given, without requiring $x\leq y$.

The final two displays give respectively the inner reserve
$n^{-1}\leq\widehat\zeta_s$, its strict inequality
$\widehat\zeta_s<\zeta_s$, the two padding inequalities, and both
output degree/regularity slack inequalities. The last two are
independent of $x$ and so hold for each $x\leq\Delta$ as required.
This checks every numerical and instance hypothesis of the local step.
Consequently, for every current touched set $T\subseteq B_s$, domain
$Z$, and actual law $\nu$ satisfying $\mathcal I(s,T;Z,\nu)$, and for
every $i\in B_s\setminus T$, that node supplies the total kernel
$\kappa$ preserving all old maps on $Z$, the conditional new-cluster
spread bound, and the new invariant for the actual law
$\nu\mathbin{\mathrm{bind}}\kappa$. This is exactly the transition
hypothesis of \noderef{ColourClassSequentialCompletion}. It does
not resample the past or condition on success for the whole class.

Put $P_s=\{j:c(j)<s\}$ and let $Z_s$ be its union of source clusters.
Since $Q\geq1$, \noderef{VariableSizeInductionBase} applies at zero:
$p_0=d$, $C/n\leq1$, the degree and colouring bounds hold, and the
given \noderef{CandidacySuperRegular} provides the original witnesses.
It gives a specified law $\nu_0$ with
$\mathcal I(0,\varnothing;Z_0,\nu_0)$, where $Z_0=\varnothing$.

Induct over $s=0,\ldots,Q-1$. Starting with a specified law
$\nu_s$ satisfying $\mathcal I(s,\varnothing;Z_s,\nu_s)$, apply
\noderef{ColourClassSequentialCompletion} to the just verified local
transition. Retain its composite kernel $K_s$ and its completed law
$\bar\nu_s=\nu_s\mathbin{\mathrm{bind}}K_s$, on the exact domain
\[
 \bar Z_s=\{x:c(B.\mathrm{sourceCluster}(x))<s
                   \text{ or }c(B.\mathrm{sourceCluster}(x))=s\}.
\]
The node gives $\mathcal I(s,B_s;\bar Z_s,\bar\nu_s)$ for that very law,
and preservation of every previously embedded map. Its empty-class
case is included: the kernel is the point-mass identity after identifying
equal domains, rather than an additional probabilistic step.
If $s+1<Q$, apply \noderef{ColourRoundReset} with
$\varepsilon_s\leq\varepsilon_{s+1}$ to this completed law.
It supplies $\mathcal I(s+1,\varnothing;\bar Z_s,\bar\nu_s)$ without
changing the domain or the law. The identity
\[
 c(j)<s\ \text{or}\ c(j)=s\quad\Longleftrightarrow\quad c(j)<s+1
\]
identifies $\bar Z_s=Z_{s+1}$, so use the same law, with only this
domain identification, as $\nu_{s+1}$. This completes the induction at
every intermediate boundary. The construction binds only the successive
local kernels, so it retains their recursively produced law throughout.

At $s=Q-1$ do not reset. Every colour value is less than $Q$, hence is
either less than $Q-1$ or equal to $Q-1$. Therefore
$\bar Z_{Q-1}=V_H$. This holds even if the last class is empty.
Identify the last law with a PMF on
$\mathcal E(V_H)=\operatorname{PartialCompatibleEmbedding}(B,A,V_H)$
without changing any underlying function. The completed invariant is
precisely $\mathcal I(Q-1,B_{Q-1};V_H,\bar\nu_{Q-1})$.
The witness clause now has no unembedded cluster, and the domain,
colouring, degree and spread clauses remain the completed invariant's
clauses. Set $\mu=\bar\nu_{Q-1}$. This proves the claimed final state
and its \noderef{VertexSpreadOn} bound. When $Q=1$ the same argument
consists of the base and the single class completion, with no reset;
no round-$Q$ invariant is ever introduced.
\end{proof}
